% =========================================================================
% capitulos/04_conclusiones_recomendaciones.tex
% Secciones Finales: Conclusiones y Recomendaciones
% =========================================================================

\section*{CONCLUSIONES}
\addcontentsline{toc}{section}{CONCLUSIONES}

\subsection*{General}
Presente una conclusión principal que dé respuesta directa al objetivo general planteado en el Capítulo III, destacando el impacto general y el valor técnico aportado a la institución receptora mediante la culminación del proyecto de prácticas preprofesionales.

\subsection*{Específicos}
Redacte una conclusión por cada objetivo específico formulado, demostrando con hechos o métricas su grado de cumplimiento:
\begin{itemize}
  \item Se levantaron y formalizaron los requerimientos técnicos y funcionales del sistema, logrando una comprensión precisa de las necesidades operativas de la entidad receptora.
  \item Se diseñó una arquitectura de software desacoplada y un modelo de datos estructurado, lo cual garantizó la consistencia de la información y la mantenibilidad del código a largo plazo.
  \item Se implementaron los módulos lógicos y servicios web planificados, permitiendo automatizar los flujos de trabajo clave y reduciendo la tasa de errores operativos.
  \item Se ejecutaron las pruebas unitarias y de integración correspondientes, verificando la robustez de los componentes antes de su puesta en marcha en el entorno de pruebas.
\end{itemize}

\clearpage
\section*{RECOMENDACIONES}
\addcontentsline{toc}{section}{RECOMENDACIONES}

Las recomendaciones deben ser propositivas, viables y orientadas a la continuidad o mejora de la labor técnica realizada:

\begin{itemize}
  \item \textbf{A la institución o empresa receptora:} Se recomienda continuar con el plan de actualización tecnológica de la infraestructura restante, así como programar capacitaciones periódicas para el personal usuario de las nuevas herramientas implementadas.
  \item \textbf{Al área técnica de desarrollo:} Se sugiere extender la cobertura de pruebas automatizadas hacia pruebas de estrés y rendimiento, garantizando un comportamiento óptimo ante picos de concurrencia en la plataforma.
  \item \textbf{A la Facultad y Escuela Profesional (UNAS):} Se recomienda seguir incentivando convenios de prácticas preprofesionales con organizaciones que promuevan la aplicación de metodologías de ingeniería y buenas prácticas arquitectónicas.
  \item \textbf{A futuros practicantes:} Se recomienda documentar continuamente los cambios y decisiones técnicas mediante repositorios versionados y estándares de código limpio, facilitando la transferencia de conocimiento entre promociones.
\end{itemize}
